\def\stat{moskin}





\def\tit{ГРАФОВЫЕ \textit{n}-ГРАММЫ В~ЗАДАЧЕ АТРИБУЦИИ ТЕКСТОВ}





\def\titkol{Графовые \textit{n}-граммы в~задаче атрибуции текстов}





\def\aut{Н.\,Д.~Москин$^1$, А.\,А.~Рогов$^2$, А.\,А.~Лебедев$^3$}





\def\autkol{Н.\,Д.~Москин, А.\,А.~Рогов, А.\,А.~Лебедев}





\titel{\tit}{\aut}{\autkol}{\titkol}





\index{Москин Н.\,Д.}


\index{Рогов А.\,А.}


\index{Лебедев А.\,А.}


\index{Moskin N.\,D.}


\index{Rogov A.\,A.}


\index{Lebedev A.\,A.}





%{\renewcommand{\thefootnote}{\fnsymbol{footnote}}


%\footnotetext[1]


%{Исследование выполнено с~использованием ЦКП <<Информатика>> ФИЦ ИУ РАН.}} 





\renewcommand{\thefootnote}{\arabic{footnote}}


\footnotetext[1]{Петрозаводский государственный университет, moskin@petrsu.ru}


\footnotetext[2]{Петрозаводский государственный университет, rogov@petrsu.ru}


\footnotetext[3]{Петрозаводский государственный университет, perevodchik88@yandex.ru}








 \label{st\stat}





 \thispagestyle{headings}


 


\vspace*{-18pt}





  





 \Abst{Представлены результаты исследований в~области моделирования структуры текстов с~использованием обобщенной контекстно-зависимой тео\-ре\-ти\-ко-гра\-фо\-вой модели. 


Объектом исследования стали в~основном литературные и~фольклорные тексты, для которых 


возникает задача атрибуции. Например, таких текстов много в~творчестве известного 


русского писателя Ф.\,М.~Достоевского. Авторы показывают, как можно построить 


гибридные модели, основанные на деревьях зависимостей, графовых моделях 


синтаксической структуры связей между простыми предложениями в~составе 


многокомпонентного сложного предложения и~графов <<силь\-ных связей>> сочетаемости 


слов различных грамматических классов. Такие модели позволяют конструировать новые 


информативные признаки, потенциально применимые в~атрибуции текстов. Примером 


служит частота встречаемости графовых $n$-грамм, которые представляют собой обобщение 


обычных $n$-грамм, синтаксических $n$-грамм и~других подобных конструкций, 


применяемых в~стилистических исследованиях.} 


 


 \KW{искусственный интеллект; атрибуция текстов; граф; метаграф; гибридный граф; 


фольклорный текст; литературный текст; графовая $n$-грамма}





\DOI{10.14357\!/\!08696527230411}{ANMCWV}  





\vskip 10pt plus 9pt minus 6pt





 \thispagestyle{headings}


 


 \vspace*{-16pt}


  


\section{Введение}





\vspace*{-2pt}








  В настоящее время анализ текстовой информации количественными 


методами особенно актуален, поскольку с~его помощью можно решать большое 


число задач разного рода, начиная от частных (среди которых машинный 


перевод, реферирование, плагиат, атрибуция, выделение неоднородных 


фрагментов в~тексте, классификация текстов и~т.\,д.), заканчивая глобальной~--- 


пониманием текста. В~отечественной лингвистике подобного рода 


исследования, проводимые на стыке филологии и~математики, имеют давнюю 


историю~[1]. 


  


  Сегодня существует большое число различных способов, позволяющих 


решать вышеуказанные частные задачи. Наиболее популярны нейросетевые 


методы. Анализируя большой объем информации, данные методы позволяют 


выявить статистические закономерности и~в~зависимости от цели задачи 


описать их или найти выбросы (отклонения от общих закономерностей). 


Возникло новое на\-прав\-ле\-ние анализа текстов~--- автоматическая обработка 


текстов (АОТ), сочетающая в~себе компьютерную (или математическую) 


лингвистику и~машинное обучение (подробнее см.~[2]).


  


  Если говорить о цели данного исследования, то в~центре внимания авторов 


находится исследование вопросов применения графовых $n$-грамм в~задаче 


атрибуции текстов. Подобные структуры позволяют сравнивать разные тексты 


на основании объективных, количественно выраженных параметров (например, 


соотношение разных частей речи в~пределах текста, длина синтаксических 


конструкций, частота связей между различными частями речи и~т.\,п.).


  





\section{Гибридные модели атрибутируемого текста}





  В отличие от других видов моделей, в~графовых моделях указываются 


взаимосвязи между конструкциями текста (словами и~предложениями). В~связи 


с~бурным развитием вычислительной техники наблюдается тенденция 


усложнения графовых структур, возникающих при моделировании сложных 


объектов на стыке математики, информационных технологий и~гуманитарных 


наук. Например, в~[1] приводятся описания таких понятий, как нечеткий граф, 


темпоральный граф и~иерархический граф. Все чаще в~научных исследованиях 


используются метаграфы~[3]: здесь метавершина может включать как вершины 


(или метавершины), так и~ребра; ребро метаграфа может соединять вершины 


внутри одной метавершины, вершины между различными метавершинами, 


метавершины, а~также вершины и~метавершины. В~[4] определяется понятие 


вложенного метаграфа ($n$-уров\-не\-во\-го графа) как обобщения графов, 


гиперграфов и~метаграфов. При этом он позволяет задавать ориентацию ребер.


  


  Рассмотренное авторами рекурсивное определение обобщенной  


кон\-текст\-но-за\-ви\-си\-мой тео\-ре\-ти\-ко-гра\-фо\-вой модели для решения 


задачи анализа текстов изложено в~[1]. Она позволяет строить новые 


гибридные модели, которые можно впоследствии использовать для получения 


новых информативных признаков для проведения атрибуции. Рассмотрим 


фрагмент текста <<Дворянин, желающий быть крестьянином>> (журнал 


<<Время>>, 1861, №\,12), который <<с~давних времен приписывается 


Достоевскому>>~\cite[с.~278]{5-mos}, однако однозначно авторство не 


установлено (см.\ работы Б.\,В.~Томашевского, В.\,С.~Нечаева, Г.~Хетсо и~др.): 


<<\textit{Антагонизмъ двухъ лагерей существуетъ; реформа 19~февраля 


стремится къ его уничтоженiю, къ сглаженiю неравноправности; но явленiе, 


укоренившееся исторiей, не сглаживается въ одинъ день или въ одинъ годъ; 


оно, старое, тогда только исчезнетъ, когда новое будетъ прiобр\hspace*{-4pt}{\raisebox{-1pt}{


\epsfxsize=1.5mm %328mm


\epsfbox{mos-yat-kos.eps}


}}\hspace*{-4pt}тено 


вполн\hspace*{-4pt}{\raisebox{-1pt}


{


\epsfxsize=1.5mm %328mm


\epsfbox{mos-yat-kos.eps}


}}}\hspace*{-4pt}>>. Число слов в~тексте $K\hm=37$. Каждому слову можно 


поставить в~соответствие часть речи: $N$~--- существительное; $A$~--- 


прилагательное; $C$~--- местоимение; $O$~--- числительное; $V$~--- глагол; 


$E$~--- причастие; $G$~--- деепричастие; $D$~--- наречие; $S$~--- категория 


состояния; $P$~--- предлог; $L$~--- союз; $U$~--- частица; $I$~--- междометие. 


В~табл.~1 представлено множество слов текста $W\hm= \{ w_k\}_{k=1}^{37}$.





\begin{table}\small %tabl1


\vspace*{-7pt}


    \begin{center}


    \Caption{Соответствие вершин и~слов текста <<\textit{Антагонизмъ двухъ 


лагерей}$\ldots$>>


\label{mos-t1-b}}


     \vspace*{2ex}


     


      \tabcolsep=6.6pt


     \begin{tabular}{|c|p{45mm}|c|c|c|}


     \hline


\textbf{№}&\multicolumn{1}{c|}{\textbf{Подмножества 


слов}}&\textbf{Вершина}&\textbf{Группы}&\tabcolsep=0pt\begin{tabular}{c}\textbf{Функция}\\ \textbf{принадлежности}\end{tabular}\\


\hline


&&&&\\[-9pt]


1&$W_1=\{w_1\}$: <<\textit{Антагонизмъ}>>&$v_1\hphantom{{}_3}$&$\hphantom{{}_3}\alpha(v_1)=N$&$\mu(v_1)=1$\\


\hline


&&&&\\[-9pt]


2&$W_2=\{w_2\}$: <<\textit{двухъ}>>&$v_2\hphantom{{}_3}$&$\hphantom{{}_3}\alpha(v_2)=O$&$\mu(v_2)=1$\\


\hline


&&&&\\[-9pt]


3&$W_3=\{w_3\}$: <<\textit{лагерей}>>&$v_3\hphantom{{}_3}$&$\hphantom{{}_3}\alpha(v_3)=N$&$\mu(v_3)=1$\\


\hline


&&&&\\[-9pt]


4&$W_4=\{w_4\}$: <<\textit{существуетъ}>>&$v_4\hphantom{{}_3}$&$\hphantom{{}_3}\alpha(v_4)=V$&$\mu(v_4)=1$\\


\hline


&&&&\\[-9pt]


$\cdots$&\multicolumn{1}{|c|}{$\cdots$}&$\cdots$&$\cdots$&$\cdots$\\


\hline


&&&&\\[-9pt]


37&$W_{37}=\{w_{37}\}$: 


<<\textit{вполн}\hspace*{-3pt}{\raisebox{-1pt}{


\epsfxsize=1.5mm


\epsfbox{mos-yat-kos.eps}


}}\hspace*{-3pt}>>&$v_{37}$&$\alpha(v_{37})=D$&$\mu(v_{37})=1$\\


\hline


&&&&\\[-9pt]


38&\multicolumn{1}{|c|}{---}&$v_{38}$&$\alpha(v_{38})=F$&$\mu(v_{38})=1$\\


\hline


&&&&\\[-9pt]


39&$W_{38}=\{w_1, w_2, w_3, w_4\}$: <<\textit{Ан\-та\-го\-низмъ}>>, <<\textit{двухъ}>>, 


<<\textit{ла\-ге\-рей}>>, <<\textit{су\-щест\-ву\-етъ}>>&$v_{39}$&---&$\mu(v_{39})=1$\\


\hline


&&&&\\[-9pt]


40&$W_{39}=\{ w_5, w_6, w_7, w_8, w_9$,\newline $w_{10}, w_{11}, w_{12}, w_{13}, w_{14}\}$: 


<<\textit{ре\-фор\-ма}>>, <<\textit{19}>>, <<\textit{фев\-ра\-ля}>>, <<\textit{стре\-мит\-ся}>>, 


<<\textit{къ}>>, <<\textit{его}>>, <<\textit{унич\-то\-же\-нiю}>>, <<\textit{къ}>>, 


<<\textit{сгла\-же\-нiю}>>, <<\textit{не\-рав\-но\-прав\-ности}>>&$v_{40}$&---


&$\mu(v_{40})=1$\\


\hline


&&&&\\[-9pt]


41&$W_{40}=\{w_{15}, w_{16}, w_{17}, w_{18}$,\newline $w_{19}, w_{20}, w_{21}, w_{22}, w_{23}, 


w_{24}$,\newline $w_{25}, w_{26}, w_{27}\}$: <<\textit{но}>>, <<\textit{явле\-нiе}>>, 


<<\textit{уко\-ре\-нив\-ше\-еся}>>, <<\textit{ис\-то\-рi\-ей}>>, <<\textit{не}>>, 


<<\textit{сгла\-жи\-ва\-ет\-ся}>>, <<\textit{въ}>>, <<\textit{одинъ}>>, <<\textit{день}>>, 


<<\textit{или}>>, <<\textit{въ}>>, <<\textit{одинъ}>>, <<\textit{годъ}>>&$v_{41}$&---


&$\mu(v_{41})=1$\\


\hline


&&&&\\[-9pt]


42&$W_{41}=\{ w_{28}, w_{29}, w_{30}, w_{31}$,\newline $w_{32}\}$: <<\textit{оно}>>, 


<<\textit{ста\-рое}>>, <<\textit{тогда}>>, <<\textit{толь\-ко}>>, 


<<\textit{ис\-чез\-нетъ}>>&$v_{42}$&---&$\mu(v_{42})=1$\\


\hline


&&&&\\[-9pt]


43&$W_{42}= \{ w_{33}, w_{34}, w_{35}, w_{36}$,\newline $w_{37}\}$: <<\textit{когда}>>, 


<<\textit{но\-вое}>>, <<\textit{бу\-детъ}>>, <<\textit{прi\-обр}\hspace*{-2.5pt}{\raisebox{-1pt}{


\epsfxsize=1.5mm


\epsfbox{mos-yat-kos.eps}


}}\hspace*{-2.5pt}\textit{те\-но}>>, 


<<\textit{впол\-н}\hspace*{-2.5pt}{\raisebox{-1pt}{


\epsfxsize=1.5mm


\epsfbox{mos-yat-kos.eps}


}}\hspace*{-2.5pt}>>&$v_{43}$&---&$\mu(v_{43})=1$\\


\hline


&&&&\\[-9pt]


44&$W_{43}= \{w_1, w_3, w_5, w_7, w_{11}$,\newline $w_{13}, w_{14}, w_{16}, w_{18}, w_{23}, 


w_{27}$,\newline


$w_{29}, w_{34}\}$: <<\textit{Ан\-та\-го\-низмъ}>>,\newline <<\textit{ла\-ге\-рей}>>, 


<<\textit{ре\-фор\-ма}>>, <<\textit{фев\-ра\-ля}>>, <<\textit{уничтоженiю}>>, 


<<\textit{сгла\-же\-нiю}>>, <<\textit{не\-рав\-но\-прав\-ности}>>, <<\textit{явле\-нiе}>>, 


<<\textit{ис\-то\-рi\-ей}>>, <<\textit{день}>>, <<\textit{годъ}>>, <<\textit{ста\-рое}>>, 


<<\textit{новое}>>&$v_{44}$&$\alpha(v_{44})=N$&$\mu(v_{44})=1$\\


\hline


\multicolumn{5}{r}{\footnotesize{\textit{Окончание табл.~$\ref{mos-t1-b}$ на с.~$\pageref{mos-t1-e}$}}}


\end{tabular}


\end{center}


\vspace*{-6pt}


\end{table}











       Первую теоретико-графовую модель построим на основе дерева 


зависимостей. Такая структура позволяет наглядно изобразить синтаксические 


связи в~предложении и~на их основе построить ряд признаков, которые можно 


использовать в~атрибуции художественных текстов~[6]. Здесь выделяются 


38 вершин $V\hm= \{v_i\}_{i=1}^{38}$, которые соответствуют 


следующим подмножествам слов~$W_i$ (см.\ табл.~1). Функция~$\alpha$ задает 


часть речи для~$v_i$. При этом буквой~$F$ обозначим фиктивную вершину, не 


связанную со словами. Функция принадлежности~$\mu$ для всех вершин 


и~ребер принимает значение~1 (объекты однозначно определены), т.\,е.\ 


    $\mu(v_i)\hm= \mu(e_j)\hm=1$.


    


    \addtocounter{table}{-1}


\begin{table}\small %tabl1


    \begin{center}


    \Caption{(\textit{окончание} Соответствие вершин и~слов текста <<\textit{Антагонизмъ двухъ 


лагерей}$\ldots$>>\label{mos-t1-e}}


     \vspace*{2ex}


     


     \tabcolsep=6.8pt


     \begin{tabular}{|c|p{45mm}|c|c|c|}


     \hline


\textbf{№}&\multicolumn{1}{c|}{\textbf{Подмножества 


слов}}&\textbf{Вершина}&\textbf{Группы}&\tabcolsep=0pt\begin{tabular}{c}\textbf{Функция}\\ \textbf{принадлежности}\end{tabular}\\


\hline


&&&&\\[-9pt]


45&$W_{44}=\{ w_2, w_6, w_{22}, w_{26}\}$: <<\textit{двухъ}>>, <<\textit{19}>>, 


<<\textit{одинъ}>>, 


<<\textit{одинъ}>>&$v_{45}$&$\alpha(v_{45})=O$&$\mu(v_{45})=1$\\


\hline


&&&&\\[-9pt]


46&$W_{45}=\{ w_{10}, w_{28}\}$: <<\textit{его}>>, 


<<\textit{оно}>>&$v_{46}$&$\alpha(v_{46})=C$&$\mu(v_{46})=1$\\


\hline


&&&&\\[-9pt]


47&$W_{46}=\{ w_4, w_8, w_{20}, w_{32},\newline w_{35}\}$: <<\textit{су\-щест\-ву\-етъ}>>, 


<<\textit{стре\-мит\-ся}>>, <<\textit{сгла\-жи\-ва\-ет\-ся}>>, <<\textit{ис\-чез\-нетъ}>>, 


<<\textit{бу\-детъ}>>&$v_{47}$&$\alpha(v_{47})=V$&$\mu(v_{47})=1$\\


\hline


&&&&\\[-9pt]


48&$W_{47}=\{ w_{17}, w_{36}\}$: <<\textit{уко\-ре\-нив\-ше\-еся}>>, 


<<\textit{прi\-об\-р}\hspace*{-2.5pt}{\raisebox{-1pt}{


\epsfxsize=1.5mm


\epsfbox{mos-yat-kos.eps}


}}\hspace*{-2.5pt}\textit{те\-но}>>&$v_{48}$&$\alpha(v_{48})=E$&$\mu(v_{48})=1$\\


\hline


&&&&\\[-9pt]


49&$W_{48}=\{w_{30}, w_{37}\}$: <<\textit{тогда}>>, 


<<\textit{впол\-н}\hspace*{-2.5pt}{\raisebox{-1pt}{


\epsfxsize=1.5mm


\epsfbox{mos-yat-kos.eps}


}}\hspace*{-2.5pt}>>&$v_{49}$&$\alpha(v_{49})=D$&$\mu(v_{49})=1$\\


\hline


&&&&\\[-9pt]


50&$W_{49}=\{ w_{19}, w_{31}\}$: <<\textit{не}>>, 


<<\textit{толь\-ко}>>&$v_{50}$&$\alpha(v_{50})=U$&$\mu(v_{50})=1$\\


\hline


&&&&\\[-9pt]


51&$W_{50} =\{ w_9, w_{12}, w_{21}, w_{25}\}$: <<\textit{къ}>>, <<\textit{къ}>>, 


<<\textit{въ}>>, <<\textit{въ}>>&$v_{51}$&$\alpha(v_{51})=P$&$\mu(v_{51})=1$\\


\hline


&&&&\\[-9pt]


52&$W_{51}=\{w_{15}, w_{24}, w_{33}\}$: <<\textit{но}>>, <<\textit{или}>>, 


<<\textit{когда}>>&$v_{52}$&$\alpha(v_{52})=L$&$\mu(v_{52})=1$\\


\hline


\end{tabular}


\end{center}


\vspace*{-12pt}


\end{table}


    


    


    Фрагмент теоретико-графовой модели предложения <<\textit{Антагонизмъ 


двухъ лагерей}$\ldots$>> представлен на рис.~1 (здесь отсутствует иерархия).





\begin{figure} %fig1


\vspace*{1pt}


  \begin{center}


 \mbox{%


 \epsfxsize=62mm 


\epsfbox{mos-1.eps}


 }


\end{center}


\vspace*{-11pt}


\Caption{Фрагмент дерева зависимости предложения <<\textit{Антагонизмъ двухъ 


лагерей}$\ldots$>>}


%\vspace*{-6pt}


%\end{figure}


%


%\begin{figure} %fig2


\vspace*{6pt}


  \begin{center}


 \mbox{%


 \epsfxsize=62mm 


\epsfbox{mos-2.eps}


 }


\end{center}


\vspace*{-11pt}


\Caption{Структурная модель сложного предложения <<\textit{Антагонизмъ двухъ 


лагерей}$\ldots$>>}


\end{figure}





\begin{table}[b]\small %tabl2


\vspace*{-18pt}


\begin{center}


\Caption{Матрица сильных связей (пороговое значение $a\hm=0{,}02$)}


\vspace*{2ex}





\tabcolsep=1.25pt %1.6pt


\begin{tabular}{|c|c|c|c|c|c|c|c|c|c|}


\hline


\tabcolsep=0pt\begin{tabular}{c}Часть\\ речи\end{tabular}&\tabcolsep=0pt\begin{tabular}{c}\textrm{Cущест-}\\ \textrm{вительное}\\ ($N$)\end{tabular}&


\tabcolsep=0pt\begin{tabular}{c}\textrm{Числи-}\\ \textrm{тельное}\\ ($O$)\end{tabular}&


\tabcolsep=0pt\begin{tabular}{c}\textrm{Место-}\\ \textrm{имение}\\ ($C$)\end{tabular}&


\tabcolsep=0pt\begin{tabular}{c}\textrm{Глагол}\\ ($V$) \end{tabular}&


\tabcolsep=0pt\begin{tabular}{c}\textrm{При-}\\ \textrm{частие}\\ ($E$)\end{tabular}&


\tabcolsep=0pt\begin{tabular}{c}\textrm{Наречие}\\ ($D$)\end{tabular}&


\tabcolsep=0pt\begin{tabular}{c}\textrm{Частица}\\ ($U$)\end{tabular}&


\tabcolsep=0pt\begin{tabular}{c}\textrm{Предлог}\\ ($P$)\end{tabular}&


\tabcolsep=0pt\begin{tabular}{c}\textrm{Союз}\\ ($L$)\end{tabular}\\


\hline


$N$&0,02778&0,05556&0,02778&0,08333&0,02778&0,02778&0,02778&0,02778&0,05556\\


$O$&0,11111&0&0&0&0&0&0&0&0\\


$C$&0,05556&0&0&0&0&0&0&0&0\\


$V$&0,02778&0&0&0&0,02778&0&0&0,05556&0,02778\\


$E$&0,02778&0&0&0&0&0,02778&0&0&0\\


$D$&0&0&0&0&0&0&0,02778&0&0\\


$U$&0&0&0&0,05556&0&0&0&0&0\\


$P$&0,02778&0,05556&0,02778&0&0&0&0&0&0\\


$L$&0,05556&0&0&0&0&0&0&0,02778&0\\


\hline


\end{tabular}


\end{center}


\end{table}





    


    Вторая модель образуется из вершин $v_{39}, \ldots , v_{43}$ и~отражает 


синтаксическую структуру связей между простыми предложениями в~составе 


многокомпонентного сложного предложения. Структурно фрагмент можно 


представить следующим образом: [Антагонизмъ двухъ лагерей 


существуетъ;]$_1$ [реформа 19~февраля стремится къ его уничтоженiю, къ 


сглаженiю неравноправности;]$_2$ [но явленiе, укоренившееся исторiей, не 


сглаживается въ одинъ день или въ одинъ годъ;]$_3$ [оно, старое, тогда только 


исчезнетъ,]$_4$ [когда новое будетъ прiобр\hspace*{-2.5pt}{\raisebox{-0.5pt}{


\epsfxsize=1.7mm


\epsfbox{mos-yat-prym.eps}


}}\hspace*{-2.5pt}тено вполн\hspace*{-2.5pt}{\raisebox{-0.5pt}{


\epsfxsize=1.7mm


\epsfbox{mos-yat-prym.eps}


}}\hspace*{-2.5pt}]$_5$. Связи 


делятся на группы: $e_{38}$~--- бессоюзные предложения (БСП) со значением 


перечисления; $e_{39}$~--- сложносочиненные предложения (ССП) 


с~противительным союзом \textit{но} (отношения сопоставления); $e_{40}$~--- 


БСП со значением пояснения; $e_{41}$~--- сложноподчиненные предложения 


(СПП) с~придаточным времени (рис.~2).


    











    Третья модель~[1] описывает, как в~тексте сочетаются слова разных 


грамматических классов (каждая вершина графа $v_{44}, v_{45}, \ldots , v_{52}$ 


соответствует определенной части речи). Если относительная частота таких 


связей превышает заранее заданный порог, то связь считается <<сильной>> 


и~отображается в~матрице (в~табл.~2 пороговое значение выбрано на уровне 


$a\hm= 0{,}02$). Количественное значение можно сохранить в~обобщенной 


модели как значение функции~$\mu$. Например, связь между 


грамматическими классами <<существительное>> и~<<чис\-ли\-тель\-ное>> 


соответствует $\mu\left( e_{43}\hm= (v_{44}, v_{45})\right)\hm= 0{,}05556$. 


Стрелки на рис.~3 отражают последовательность встречаемости двух слов 


определенных частей речи (биграмму).


    














\begin{figure} %fig3


\vspace*{1pt}


  \begin{center}


 \mbox{%


 \epsfxsize=61.793mm 


\epsfbox{mos-3.eps}


 }


\end{center}


\vspace*{-11pt}


\Caption{Граф <<сильных>> связей}


\vspace*{-9pt}


\end{figure}





%\vspace*{-3pt}





\section{Графовые {\boldmath{$n$}}-граммы}





%\vspace*{-3pt}





  К одним из значимых характеристик, используемых при решении задачи 


атрибуции текстов, относятся $n$-грам\-мы (см. примеры  


в~\cite[с.~152--154]{5-mos}). Это последовательности текстовых элементов 


(букв, слов, лемм, тегов частей речи и~т.\,д.), взятых в~порядке их появления 


в~тексте. Значение~$n$ указывает, сколько элементов должно быть взято, 


и~определяет размер последовательности. В~терминах обобщенной модели 


можно определить $n$-грам\-му сле\-ду\-ющим образом: 


  \begin{enumerate}[(1)]


  \item  первый случай: биграмма (2-грамма). Определим две части речи~$t_1$ 


и~$t_2$. Выделим два подмножества слов текста~$W_1^\prime$ 


и~$W_2^\prime$, которые соответствуют частям речи~$t_1$ и~$t_2$, т.\,е.\ для 


любых $w_i\hm\in W_1^\prime$ и~$w_j\hm\in W_2^\prime$ выполняется 


$\alpha(w_i)\hm=t_1$ и~$\alpha(w_j)\hm=t_2$. Если слова~$w_i$ и~$w_j$ идут 


последовательно, т.\,е.\  $j\hm= i\hm+1$, то они образуют биграмму. Кстати, 


части речи и~соответственно подмножества слов~$W_1^\prime$ 


и~$W_2^\prime$ могут совпадать; 


  \item  второй случай: триграмма (3-грамма). На этот раз определим три части 


речи: $t_1$, $t_2$ и~$t_3$. Выделим три подмножества слов 


текста~$W_1^\prime$, $W_2^\prime$ и~$W_3^\prime$, которые соответствуют 


частям речи $t_1$, $t_2$ и~$t_3$, т.\,е.\ для любых $w_i\hm\in W_1^\prime$, 


$w_j\hm\in W_2^\prime$ и~$w_k\hm\in W^\prime_3$ выполняется 


$\alpha(w_i)\hm=t_1$, $\alpha(w_j)\hm=t_2$ и $\alpha(w_k)\hm=t_3$. Если 


слова~$w_i$, $w_j$ и~$w_k$ идут последовательно, т.\,е.\ $k\hm= j\hm+1\hm= 


i\hm+2$, то они образуют триграмму;


  \item  третий (общий) случай: $n$-грамма. Рассмотрим $n$ частей речи: $t_1, 


t_2, \ldots , t_n$. Выделим $n$ подмножеств слов текста $W_1^\prime, 


W_2^\prime, \ldots , W_n^\prime$, которые соответствуют частям речи $t_1, t_2, 


\ldots , t_n$, т.\,е.\ для любых $w_{ij}\hm\in W_i^\prime$ выполняется 


$\alpha(w_{ij})\hm= t_i$. Если $n$ слов соответственно из $W_1^\prime, 


W_2^\prime, \ldots , W_n^\prime$ идут последовательно, то они образуют $n$-грамму.


  \end{enumerate}


  


 


\begin{floatingfigure}{56mm} %fig4


\vspace*{1pt}


  \begin{center}


 \mbox{%


 \epsfxsize=53.631mm 


\epsfbox{mos-4.eps}


 }


\end{center}


\vspace*{-11pt}


\Caption{Графовая $n$-грамма (первый пример)}


\vspace*{9pt}


\end{floatingfigure}    


  Традиционные $n$-граммы строятся на синтагматических связях 


(последовательностях слов в~тексте), а~потому не позволяют в~полной мере 


учитывать взаимосвязи между словами, связанными подчинительной связью, 


но при этом удаленными друг от друга (это особенно важно для текстов, 


написанных на русском языке, поскольку для них характерен достаточно 


свободный порядок слов). 


Назовем \textit{графовой\linebreak $n$-грам\-мой} связный 


подграф обобщенной тео\-ре\-ти\-ко-гра\-фо\-вой модели, со\-сто\-ящий 


из~$n$~вер\-шин, каждой из которых ставится в~соответствие подмножество 


слов в~тексте. Рассмотрим пример: упорядоченному подмножеству слов 


<<\textit{реформа 19~февраля стремится}>> будет соответствовать графовая 


$n$-грам\-ма, изображенная на рис.~4 (при этом традиционная 4-грам\-ма 


выглядит как  


<<су\-щест\-ви\-тель\-ное--чис\-ли\-тель\-ное--су\-щест\-ви\-тель\-ное--гла\-гол>>, т.\,е.\ строка  


$N$--$O$--$N$--$V$ без учета связей).


 





  Если задача поиска традиционной $n$-грам\-мы сводилась к~поиску 


подстроки в~строке, то здесь мы имеем дело с~поиском изоморфного подграфа, 


в~котором вершины упорядочены. Отметим, что подобные модели 


в~упрощенном виде ранее уже применялись в~стилистических исследованиях. 


Например, в~[7] автор рассматривает такие <<конструкции>>, которые 


записываются <<$\ldots$в~виде двоеточия, слева от которого указывается 


число членов, расположенных перед сказуемым, а~после двоеточия~--- после 


него>>, а~затем показывает, как подсчитываются их частоты и~строится 


ранговое распределение линеаризованных синтаксических конструкций. 


Другие примеры можно найти в~работах Г.\,О.~Сидорова~[8], где автор 


описывает синтаксические $n$-грам\-мы, построенные нелинейным способом 


на основе анализа путей в~синтаксических деревьях. Нелинейный принцип 


построения лежит и~в~основе так называемых \textit{skipgram}~[9], где 


в~последовательности элементов случайным образом пропускаются некоторые 


из них. 


  


  Рассмотрим, какие преимущества дают гибридные модели для построения 


различных видов графовых $n$-грамм. Вернемся в~примеру на рис.~4. Как 


видно по табл.~1, словам $w_5$, $w_6$, $w_7$ и $w_8$ соответствуют не только 


вершины $v_5$, $v_6$, $v_7$ и $v_8$, но и~вершина~$v_{40}$ из второй модели, 


объединяющая данные слова. Таким образом, можно определить два условия 


формирования графовой $n$-граммы:


  \begin{enumerate}[(1)]


  \item слова $w_5$, $w_6$, $w_7$ и $w_8$ образуют графовую $n$-грам\-му тогда 


и~только тогда, когда существует некоторое подмножество~$W_i$, 


соответствующее вершине~$v_j$ второй модели, куда входят эти слова. Более 


строгое условие заключается в~том, что слова образуют данное 


подмножество~$W_i$, т.\,е.\ $W_i\hm= \{ w_5, w_6, w_7, w_8\}$, однако для 


данного примера оно не выполняется;%\\[-18pt]


  \item слова $w_5$, $w_6$, $w_7$ и $w_8$ образуют графовую $n$-грам\-му тогда 


и~только тогда, когда не существует такого подмножества~$W_i$, 


соответствующего вершине~$v_j$ второй модели, куда бы входили все эти 


слова. Например, для подстроки <<\textit{только $(w_{31})$ исчезнетъ 


$(w_{32})$, когда $(w_{33})$ новое $(w_{34})$ будетъ $(w_{35})$ прiобр\hspace*{-3pt}{\raisebox{-0.5pt}{


\epsfxsize=1.7mm


\epsfbox{mos-yat-kos.eps}


}}\hspace*{-3pt}тено} 


($w_{36}$)>> (U--V--L--N--V--E) слова $w_{31}, w_{32}\hm\in W_{41}$, тогда 


как $w_{33}, w_{34}, w_{35}, w_{36}\hm\in W_{42}$ (рис.~5).%\\[-13pt]


  \end{enumerate}


  


   \begin{floatingfigure}{64.9mm} %fig5


\vspace*{1pt}


  \begin{center}


 \mbox{%


 \epsfxsize=61.653mm 


\epsfbox{mos-5.eps}


 }


\end{center}


\vspace*{-11pt}


\Caption{Графовая $n$-грамма (второй \mbox{пример})}


\vspace*{3pt}


\end{floatingfigure}


    С другой стороны, если обратиться к~третьей модели, слова $w_5$, $w_6$, $w_7$ и~$w_8$ также соответствуют вершинам $v_{44}$, $v_{45}$ и~$v_{47}$. Между этими 


вершинами существуют <<сильные связи>>: между существительными~--- 


0,02778, между числительным и~существительным~--- 0,11111, между 


существительным и~глаголом~--- 0,02778 (см.\ табл.~2). Поэтому такую 


графовую %\linebreak 


$n$-грам\-му можно использовать в~исследовании. Если же на 


рисунке обнаруживается <<слабая связь>> ($a\hm< 0{,}02$), то графовая  


$n$-грам\-ма уже далее не анализируется. Например, на рис.~5 отсутствует 


<<сильная связь>> между союзом и~глаголом (значение равно~0).


  








  Подобный способ определения $n$-грамм потенциально представляется 


более чувствительным к~изменению авторского стиля, так как помимо 


синтагматических связей учитывает ряд других факторов, связанных 


с~построением предложений, от стремления сочетать подчинительной связью 


те или иные части речи до предрасположенности к~инверсии~--- сознательной 


перестановке слов или словосочетаний~--- как стилистическому приему, 


отличающему индивидуально-авторский стиль. Таким образом, можно 


получить более совершенный инструмент для идентификации автора, 


поскольку отличия в~ин\-ди\-ви\-ду\-аль\-но-ав\-тор\-ском стиле могут 


проявляться не только в~определенных последовательностях частей речи, но 


и~в~изменении порядка слов, а~также в~особом расположении слов 


в~предложении (в качестве примера можно сравнить общеизвестную строку 


с~инверсией <<\textit{Я~памятник себе воздвиг нерукотворный}>> из 


стихотворения А.\,С.~Пушкина и~нормативное <<\textit{Я~воздвиг себе 


нерукотворный памятник}>>). Таким образом, подобный анализ становится 


особенно актуален для поэтических текстов, в~которых инверсия~--- весьма 


распространенное явление.


  


  \vspace*{-9.5pt}


  


\section{Заключение}





\vspace*{-7pt}





  Значимое практическое преимущество рассмотренной графовой модели 


заключается в~возможности синтеза различных подходов к~анализируемым 


текстам. В~частности, в~поле зрения исследователя оказываются 


и~морфологические параметры (распределение слов текста по частям речи), 


и~синтаксические структуры, указывающие на взаимосвязь слов в~пределах 


предложения. Комбинация двух этих составляющих грамматического уровня 


языка позволяет отыскивать как типичные для определенной эпохи или стиля 


черты, так и~уникальные со\-став\-ля\-ющие, отличающие ин\-ди\-ви\-ду\-аль\-но-ав\-тор\-ский стиль того или иного писателя или поэта, что позволяет более 


эффективно решать вопросы, связанные с~атрибуцией текстов. Проведенные 


эксперименты на большом объеме текстового материала~\cite{1-mos, 5-mos} 


подтверждают выводы авторов.


  


  Таким образом, исследование фольклорных и~литературных текстов 


с~использованием обобщенной кон\-текст\-но-за\-ви\-си\-мой  


тео\-ре\-ти\-ко-гра\-фо\-вой модели позволяет строить гиб\-рид\-ные структуры, 


основанные на различных тео\-ре\-ти\-ко-гра\-фо\-вых моделях языковой 


структуры, что позволяет извлекать новые знания о~стилистической структуре. 


Одним из потенциально эффективных инструментов представляется 


извлечение и~подсчет встречаемости в~текстовых фрагментах графовых  


$n$-грамм, которые являются обобщением традиционных $n$-грамм, 


синтаксических $n$-грамм и~других подобных конструкций. Также рассмотрен 


формат для хранения текстов, их обобщенных графовых моделей и~графовых  


$n$-грамм. Данный инструмент предлагает верифицируемые, опирающиеся на 


объективные грамматические параметры текста данные, которые в~дальнейшем 


могут быть использованы в~математико-лингвистических исследованиях. 


В~част\-ности, идеи, изложенные в~данной работе, можно использовать при 


решении различных задач анализа текстов (например, поиск плагиата, 


неоднородных текстовых фрагментов, атрибуции, анализа тональности текстов 


и~др.).


  


{\small\frenchspacing


{%\baselineskip=10.8pt


\addcontentsline{toc}{section}{Литература}


\begin{thebibliography}{9}





\bibitem{1-mos}


\Au{Москин Н.\,Д.} Теоретико-графовые модели, методы и~программные средства 


интеллектуального анализа текс\-то\-вой информации на примере фольклорных и~литературных 


произведений: Дис.\ \ldots\  д-ра техн. наук.~--- Петрозаводск, 2022. 370~с.


\bibitem{2-mos}


\Au{Белов С.\,Д., Зрелова Д.\,П., Зрелов П.\,В., Кореньков~В.\,В.} Обзор методов 


автоматической обработки текс\-тов на естественном языке~/\!\!/ Системный анализ в~науке 


и~образовании, 2020. №\,3. С.~8--22. doi: 10.37005\!/\!2071-9612-2020-3-8-22. EDN: YJFAYK.





\bibitem{3-mos}


\Au{Целых А.\,А., Дедюлина М.\,А.} Тео\-ре\-ти\-ко-гра\-фо\-вые подходы к~моделированию  


актор-се\-тей в~исследованиях науки и~технологий~/\!\!/ Моделирование, оптимизация 


и~информационные технологии, 2018. Т.~6. №\,4. С.~244--259. doi:  


10.26102\!/\!2310-6018/2018.23.4.019. EDN: YZSOGL.      





\bibitem{4-mos}


\Au{Астанин С.\,В., Драгныш Н.\,В., Жуковская~Н.\,К.} Вложенные метаграфы как модели 


сложных объектов~/\!\!/ Инженерный вестник Дона, 2012. №\,4. Ч.~2. Ст.~1434. 5~с. EDN: 


PVJCYN. 





\bibitem{5-mos}


\Au{Рогов А.\,А., Абрамов Р.\,В., Бучнева Д.\,Д., Захарова~О.\,В., Ку\-ла\-ков~К.\,А., 


Лебедев~А.\,А., Мос\-кин~Н.\,Д., От\-ли\-ван\-чик~А.\,В., Са\-ви\-нов~Е.\,Д., Си\-до\-ров~Ю.\,В.}  


Проб\-ле\-ма атрибуции в~журналах <<Время>>, <<Эпоха>> и~еженедельнике 


<<Гражданин>>.~--- Петрозаводск: Острова, 2021. 391~с. EDN: AIACMP.





\bibitem{6-mos}


\Au{Севбо И.\,П.} Графическое представление синтаксических структур и~стилистическая 


диагностика.~--- Киев: Наукова думка, 1981. 192~с.


\bibitem{7-mos}


\Au{Мартыненко Г.\,Я.} Методы математической лингвистики в~стилистических 


исследованиях.~--- СПб.: Нестор-История, 2019. 296~с.


\bibitem{8-mos}


\Au{Сидоров Г.\,О.} Синтаксические $n$-грам\-мы в~компьютерной лингвистике.~--- М.: Изд-во 


Московского ун-та, 2018. 120~с.


\bibitem{9-mos}


\Au{Cheng W., Greaves~C., Warren~M.} From $n$-gram to skipgram to concgram~/\!\!/ Int. J. 


Corpus Linguis., 2006. Vol.~11. Iss.~4. P.~411--433. doi: 10.1075\!/\!ijcl.11.4.04che.





   \end{thebibliography}


} }











\vspace*{-6pt}





\hfill{\small\textit{Поступила в~редакцию 01.07.23}}





\vspace*{8pt}





\hrule





\vspace*{2pt}











\hrule





\vspace*{8pt}





%\newpage





%\vspace*{-28pt}





\def\tit{GRAPH \textit{n}-GRAMS IN~THE~TEXT ATTRIBUTION PROBLEM}











\def\titkol{Graph \textit{n}-grams in~the~text attribution problem}





\def\aut{N.\,D.~Moskin, A.\,A.~Rogov, and~A.\,A.~Lebedev}





\def\autkol{N.\,D.~Moskin, A.\,A.~Rogov, and~A.\,A.~Lebedev}





\titel{\tit}{\aut}{\autkol}{\titkol}





\vspace*{-8pt}





\noindent


Petrozavodsk State University, 33~Lenina Prosp., Petrozavodsk 185910, Russian Federation








\def\leftfootline{\small{\textbf{\thepage}


\hfill Sistemy i Sredstva Informatiki~---


Systems and Means of Informatics 2023 vol~33 no\,4}


}%


 \def\rightfootline{\small{Sistemy i Sredstva Informatiki~---


 Systems and Means of Informatics 2023 vol~33 no\,4


\hfill \textbf{\thepage}}}





\vspace*{3pt}








  





\Abste{The paper presents the results of research in the field of modeling the structure of texts 


using a generalized context-dependent graph-theoretic model. The object of the study is mainly 


literary and folklore texts for which the task of attribution arises. For example, there are many such 


texts in the works of the famous Russian writer F.\,M.~Dostoevsky. The authors show how it is 


possible to build hybrid models based on dependency trees, graph models of syntactic structure 


of links between simple sentences in a multicomponent complex sentence, and ``strong links'' graphs 


of word combinations of different grammatical classes. Such models make it possible to construct 


new informative features that are potentially applicable in the attribution of texts. An example is the 


frequency of occurrence of graph $n$-grams which are generalizations of ordinary $n$-grams syntactic 


$n$-grams, and other similar constructions used in stylistic studies. The article also discusses the 


format for storing texts, their generalized graph models, and graph $n$-grams.}





\KWE{artificial intelligence; text attribution; graph; metagraph; hybrid graph; folklore text; literary 


text; graph $n$-gram}





\DOI{10.14357\!/\!08696527230411}{ANMCWV}  





%\vspace*{-10pt}








%\Ack


%\noindent





%\vspace*{-3pt}





\renewcommand{\bibname}{\protect\rmfamily References}


%\renewcommand{\bibname}{\large\protect\rm References}





{\small\frenchspacing


{ %\baselineskip=12pt


\addcontentsline{toc}{section}{References}


\begin{thebibliography}{9}


\bibitem{1-mos-1}


\Aue{Moskin, N.\,D.} 2022. Teoretiko-grafovye modeli, metody i~programmnye sredstva 


intellektual'nogo analiza tekstovoy informatsii na primere fol'klornykh i~literaturnykh proizvedeniy 


[Graph-theoretical models, methods, and software tools for intellectual analysis of textual 


information on the example of folklore and literary works]. Petrozavodsk. D.Sc. Diss. 370~p.    


\bibitem{2-mos-1}


\Aue{Belov, S.\,D., D.\,P.~Zrelova, P.\,V.~Zrelov, and V.\,V.~Ko\-ren'\-kov.} 2020. Obzor 


metodov avtomaticheskoy obrabotki tekstov na estestvennom yazyke [Overview of methods for 


automatic natural language text processing]. \textit{Sistemnyy analiz v~nauke i~obrazovanii} 


[System Analysis in Science and Education] 3:8--22. doi: 10.37005\!/\!2071-9612-2020-3-8-22. 


EDN: YJFAYK.


\bibitem{3-mos-1}


\Aue{Tselykh, A.\,A., and M.\,A.~Dedulina}. 2018. Teoretiko-grafovye podkhody 


k~mo\-de\-li\-ro\-va\-niyu aktor-setey v~issledovaniyakh nauki i~tekhnologiy [Graph-theoretic approaches 


to modeling actor-networks in science and technology studies]. \textit{Modelirovanie, op\-ti\-mi\-za\-tsiya 


i~informatsionnye tekhnologii} [Modeling, Optimization and Information Technology]  


6(4):244--259. doi: 10.26102\!/\!2310-6018/2018.23.4.019. EDN: \mbox{YZSOGL}.      


\bibitem{4-mos-1}


\Aue{Astanin, S.\,V., N.\,V.~Dragnysh, and N.\,K.~Zhukovskaya.} 2012. Vlozhennye metagrafy 


kak modeli slozhnykh ob''ektov [Nested metagraphs as models of complex object]. 


\textit{Inzhenernyy vestnik Dona} [Engineering J. of Don] 4-2:1434. 5~p. EDN: PVJCYN. 


\bibitem{5-mos-1}


\Aue{Rogov, A.\,A., R.\,V.~Abramov, D.\,D.~Buchneva, O.\,V.~Zakharova, K.\,A.~Kulakov, 


A.\,A.~Lebedev, N.\,D.~Moskin, A.\,V.~Otlivanchik, E.\,D.~Savinov, and Yu.\,V.~Sidorov}. 2021. 


\textit{Problema atributsii v~zhur\-na\-lakh ``Vremya,'' ``Epo\-kha'' i~ezhenedel'nike ``Grazhdanin''} [The 


problem of attribution in the magazines ``Time,'' ``Epoch,'' and the weekly ``Citizen'']. 


Petrozavodsk: Ostrova. 391~p. EDN: AIACMP.


\bibitem{6-mos-1}


\Aue{Sevbo, I.\,P.} 1981. \textit{Graficheskoe predstavlenie sintaksicheskikh struktur 


i~stilisticheskaya diagnostika} [Graphic representation of syntactic structures and stylistic 


diagnostics]. Kyiv: Naukova dumka. 192~p.


\bibitem{7-mos-1}


\Aue{Martynenko, G.\,Ya.} 2019. \textit{Metody matematicheskoy lingvistiki v~stilisticheskikh 


issledovaniyakh} [Methods of mathematical linguistics in stylistic research]. Saint Petersburg: 


Nestor-History Publishing House. 296~p.


\bibitem{8-mos-1}


\Aue{Sidorov, G.\,O.} 2018. \textit{Sintaksicheskie n-grammy v~komp'yuternoy lingvistike} 


[Syntactic $n$-grams in computational linguistics]. Moscow: Moscow University Press. 120~p.


\bibitem{9-mos-1}


\Aue{Cheng, W., C.~Greaves, and M.~Warren.} 2006. From $n$-gram to skipgram to concgram. 


\textit{Int. J. Corpus Linguis.} 11(4):411--433. doi: 10.1075\!/\!ijcl.11.4.04che.





\end{thebibliography}


} }





\vspace*{-8pt}





\hfill{\small\textit{Received July 1, 2023}} 





\vspace*{-16pt} 


  


\Contr





\vspace*{-3pt}





\noindent


\textbf{Moskin Nikolai D.} (b.\ 1980)~--- Doctor of Science in technology, associate professor, 


Petrozavodsk State University, 33~Lenina Prosp., Petrozavodsk 185910, Russian Federation; 


\mbox{moskin@petrsu.ru}





\vspace*{3pt}





\noindent


\textbf{Rogov Alexander A.} (b.\ 1959)~--- Doctor of Science in technology, head of department, 


Petrozavodsk State University, 33~Lenina Prosp., Petrozavodsk 185910, Russian Federation; 


\mbox{rogov@petrsu.ru}





\vspace*{3pt}





\noindent


\textbf{Lebedev Alexander A.} (b.\ 1988)~--- Candidate of Science (PhD) in philology, associate 


professor, Petrozavodsk State University, 33~Lenina Prosp., Petrozavodsk 185910, Russian Federation; 


\mbox{perevodchik88@yandex.ru}








\label{end\stat}





\renewcommand{\bibname}{\protect\rm Литература}


